\documentclass[10pt,a4paper]{amsart}
\usepackage{amsmath,amssymb,amsthm,graphicx}
% Retained minimal preamble for the Stacks Project chapter more-algebra.tex (More on Algebra).
% The chapter file pinned-original-source.tex carries only \input{preamble} on its first line, so no
% preamble text lives inside the chapter itself.  The Stacks Project uses one single shared preamble for
% every chapter; the snapshot stages it identically as sources/stacks-chow/preamble.tex and
% sources/stacks-etale-cohomology/preamble.tex.  The block below is copied mechanically, line for line,
% from sources/stacks-chow/preamble.tex lines 14-18 (the Xy-pic setup) and 154-217 (the native theorem
% declarations together with the mathematical macros).  Nothing was added, removed or redefined; the
% wrapper class and the amsmath/amssymb/amsthm/graphicx base packages are supplied by the assembler.
% Chapter: more-algebra.tex -- More on Algebra.
% For commutative diagrams we use Xy-pic
\usepackage[all]{xy}

% We use 2cell for 2-commutative diagrams.
\xyoption{2cell}

% Theorem environments.
%
\theoremstyle{plain}
\newtheorem{theorem}[subsection]{Theorem}
\newtheorem{proposition}[subsection]{Proposition}
\newtheorem{lemma}[subsection]{Lemma}

\theoremstyle{definition}
\newtheorem{definition}[subsection]{Definition}
\newtheorem{example}[subsection]{Example}
\newtheorem{exercise}[subsection]{Exercise}
\newtheorem{situation}[subsection]{Situation}

\theoremstyle{remark}
\newtheorem{remark}[subsection]{Remark}
\newtheorem{remarks}[subsection]{Remarks}

\numberwithin{equation}{subsection}

% Macros
%
\def\lim{\mathop{\mathrm{lim}}\nolimits}
\def\colim{\mathop{\mathrm{colim}}\nolimits}
\def\Spec{\mathop{\mathrm{Spec}}}
\def\Hom{\mathop{\mathrm{Hom}}\nolimits}
\def\Ext{\mathop{\mathrm{Ext}}\nolimits}
\def\SheafHom{\mathop{\mathcal{H}\!\mathit{om}}\nolimits}
\def\SheafExt{\mathop{\mathcal{E}\!\mathit{xt}}\nolimits}
\def\Sch{\mathit{Sch}}
\def\Mor{\mathop{\mathrm{Mor}}\nolimits}
\def\Ob{\mathop{\mathrm{Ob}}\nolimits}
\def\Sh{\mathop{\mathit{Sh}}\nolimits}
\def\NL{\mathop{N\!L}\nolimits}
\def\CH{\mathop{\mathrm{CH}}\nolimits}
\def\proetale{{pro\text{-}\acute{e}tale}}
\def\etale{{\acute{e}tale}}
\def\QCoh{\mathit{QCoh}}
\def\Ker{\mathop{\mathrm{Ker}}}
\def\Im{\mathop{\mathrm{Im}}}
\def\Coker{\mathop{\mathrm{Coker}}}
\def\Coim{\mathop{\mathrm{Coim}}}

% Boxtimes
%
\DeclareMathSymbol{\boxtimes}{\mathbin}{AMSa}{"02}

%
% Macros for moduli stacks/spaces
%
\def\QCohstack{\mathcal{QC}\!\mathit{oh}}
\def\Cohstack{\mathcal{C}\!\mathit{oh}}
\def\Spacesstack{\mathcal{S}\!\mathit{paces}}
\def\Quotfunctor{\mathrm{Quot}}
\def\Hilbfunctor{\mathrm{Hilb}}
\def\Curvesstack{\mathcal{C}\!\mathit{urves}}
\def\Polarizedstack{\mathcal{P}\!\mathit{olarized}}
\def\Complexesstack{\mathcal{C}\!\mathit{omplexes}}
% \Pic is the operator that assigns to X its picard group, usage \Pic(X)
% \Picardstack_{X/B} denotes the Picard stack of X over B
% \Picardfunctor_{X/B} denotes the Picard functor of X over B
\def\Pic{\mathop{\mathrm{Pic}}\nolimits}
\def\Picardstack{\mathcal{P}\!\mathit{ic}}
\def\Picardfunctor{\mathrm{Pic}}
\def\Deformationcategory{\mathcal{D}\!\mathit{ef}}

% --- appended 2026-10-12 (Chief): the project own definitions of the reference/slogan/history
% display environments, copied mechanically from the shared Stacks preamble lines 8-12
% (\usepackage{verbatim} plus the three \newenvironment lines, which the chapter preamble keeps
% verbatim from the project preamble). Nothing added, removed or redefined: these environments
% carry no mathematical content and the project itself typesets them as comments. Without these
% three lines a counted statement whose own text uses them fails with Environment slogan undefined.
\usepackage{verbatim}
\newenvironment{reference}{\comment}{\endcomment}
\newenvironment{slogan}{\comment}{\endcomment}
\newenvironment{history}{\comment}{\endcomment}
\usepackage{cleveref}
\crefname{originallemma}{original lemma}{original lemmas}
\Crefname{originallemma}{Original lemma}{Original lemmas}
\makeatletter
\newlabel{lemma-sequence-Koszul-complexes}{{0921}{source}{}{sourcepointer-1}{}}
\newlabel{lemma-sequence-Koszul-complexes@cref}{{[originallemma][0][]0921}{[1][0][]source}{}{}{}}
\makeatother
\crefname{originallemma}{original lemma}{original lemmas}
\Crefname{originallemma}{Original lemma}{Original lemmas}
\makeatletter
\newlabel{derived-lemma-pro-isomorphism-bis}{{0G3D}{source}{}{sourcepointer-2}{}}
\newlabel{derived-lemma-pro-isomorphism-bis@cref}{{[originallemma][0][]0G3D}{[1][0][]source}{}{}{}}
\makeatother
\crefname{originallemma}{original lemma}{original lemmas}
\Crefname{originallemma}{Original lemma}{Original lemmas}
\makeatletter
\newlabel{lemma-consequence-Artin-Rees}{{0EGU}{source}{}{sourcepointer-3}{}}
\newlabel{lemma-consequence-Artin-Rees@cref}{{[originallemma][0][]0EGU}{[1][0][]source}{}{}{}}
\makeatother
\title{The derived tensor product with the quotients by the powers of an ideal agrees with the plain quotients as a pro-system}
\author{The Stacks Project Authors}
\date{stacks-project:more-algebra.tex:local-snapshot-20261012; GNU FDL 1.2 or later}
\begin{document}
\maketitle
\setcounter{subsection}{1}
\noindent\textit{Editorial note.} All delivered mathematical text is the author's own, copied from the pinned Stacks Project chapter file more-algebra.tex; no statement or proof step was rewritten, reordered or removed. The counted statement is the labeled lemma environment at source lines 29014-29023 (canonical tag 0G9M) and the counted proof is the proof environment at source lines 29025-29057, which immediately follows it and is its own proof.\par\medskip
\noindent\textit{Editorial note.} The preamble span runs to line 90, the current end of the retained preamble, which carries the shared Stacks preamble together with the appended definitions of the reference, slogan and history environments; the counted statement and proof do not themselves use those environments, but the span is taken whole so that the file version used for this card is unambiguous.\par\medskip
\noindent\textit{Editorial note.} The proof refers three times, twice to results of this chapter and once to a criterion of another chapter; all three are navigation stubs with their roles recorded in the reasons.\par\medskip
\noindent\textit{Editorial note.} The source URL is the source-level Stacks Project URL, as required when more than one card is produced from one source; the per-claim tag is carried by claim\_key and by this note. The counted statement and proof cite no bibliography entry, so this card needs no bibliography.\par\medskip
\noindent\textit{Allowed standard prerequisites (editorial).}\begin{itemize}
\item The inverse system of quotients by the powers of an ideal is pro-isomorphic to the inverse system of Koszul complexes on the powers of a generating set. invoked at proof lines 29026-29031. Replaces the quotient system by the Koszul system, which is a complex of finite free modules sitting in finitely many degrees.
\item The Koszul complex on the n-th powers of a set of r generators sits in degrees from -r to 0, so source and target of the comparison have cohomology concentrated in a fixed bounded range. used at proof lines 29032-29037. Makes the pro-isomorphism criterion applicable by bounding the degrees in which the systems live.
\item An isomorphism of pro-objects of a derived category is detected on cohomology in each degree, and one may replace a bounded complex by a bounded above complex of finite free modules. used at proof lines 29037-29045. Reduces the statement to a pro-isomorphism of modules in each degree for a convenient representative of the complex.
\end{itemize}\medskip
\ifdefined\hypertarget\hypertarget{sourcepointer-1}{}\fi
\noindent\textit{Source reference.} Original lemma 0921: sources/stacks-more-algebra/pinned-original-source.tex, source lines 27172--27172. The reference is invoked at proof lines 29029-29031 to replace the system of quotients by the Koszul system; every later step of the proof uses the replacement and stays inside the counted span.\par\medskip
\ifdefined\hypertarget\hypertarget{sourcepointer-2}{}\fi
\noindent\textit{Source reference.} Original lemma 0G3D: sources/stacks-chow/tags.txt, source lines 20088--20088. The reference is the criterion that lets a pro-isomorphism of complexes be checked on cohomology in each degree, applied at proof lines 29037-29038; the degree bound that makes it applicable is established inside the counted span.\par\medskip
\ifdefined\hypertarget\hypertarget{sourcepointer-3}{}\fi
\noindent\textit{Source reference.} Original lemma 0EGU: sources/stacks-more-algebra/pinned-original-source.tex, source lines 28867--28867. The reference is applied at proof lines 29051-29056 to identify both sides of the comparison in each degree with the pro-system of quotients of the cohomology; the reductions that bring the statement into that form are inside the counted span.\par\medskip
\begin{lemma}
\label{lemma-derived-completion-plain-completion}
Let $A$ be a Noetherian ring. Let $I \subset A$ be an ideal. Let
$M^\bullet$ be a bounded complex of finite $A$-modules. The
inverse system of maps
$$
M^\bullet \otimes_A^\mathbf{L} A/I^n \longrightarrow M^\bullet/I^nM^\bullet
$$
defines an isomorphism of pro-objects of $D(A)$.
\end{lemma}

\begin{proof}
Say $I = (f_1, \ldots, f_r)$.
Let $K_n \in D(A)$ be the object represented by the Koszul complex
on $f_1^n, \ldots, f_r^n$. Recall that we have maps $K_n \to A/I^n$
which induce a pro-isomorphism of inverse systems, see
Lemma \ref{lemma-sequence-Koszul-complexes}.
Hence it suffices to show that
$$
M^\bullet \otimes_A^\mathbf{L} K_n \longrightarrow M^\bullet/I^nM^\bullet
$$
defines an isomorphism of pro-objects of $D(A)$. Since $K_n$ is represented
by a complex of finite free $A$-modules sitting in degrees $-r, \ldots, 0$
there exist $a, b \in \mathbf{Z}$ such that
the source and target of the displayed arrow have vanishing cohomology
in degrees outside $[a, b]$ for all $n$. Thus we may apply
Derived Categories, Lemma \ref{derived-lemma-pro-isomorphism-bis}
and we find that it suffices to show that the maps
$$
H^i(M^\bullet \otimes_A^\mathbf{L} A/I^n) \to H^i(M^\bullet/I^nM^\bullet)
$$
define isomorphisms of pro-systems of $A$-modules for any $i \in \mathbf{Z}$.
To see this choose a quasi-isomorphism
$P^\bullet \to M^\bullet$ where $P^\bullet$ is a bounded
above complex of finite free $A$-modules. The arrows above are
given by the maps
$$
H^i(P^\bullet/I^nP^\bullet) \to H^i(M^\bullet/I^nM^\bullet)
$$
These define an isomorphism of pro-systems by
Lemma \ref{lemma-consequence-Artin-Rees}. Namely,
the lemma shows both are isomorphic to the pro-system $H^i/I^nH^i$
with $H^i = H^i(M^\bullet) = H^i(P^\bullet)$.
\end{proof}

\end{document}
